% Sommario (abstract) dell'allegato: solo il testo, da racchiudere in \begin{abstract}...\end{abstract} o in un paragrafo «Sommario».
Si descrive un prototipo di percezione che rileva elementi mobili e ostacoli con una camera di profondità a stereo attivo, frontale e ferma, e che continua a dare un'indicazione dichiarata quando un ostacolo nasconde il soggetto. Il sistema tiene distinte tre informazioni: l'osservato; il non osservato ma previsto, con una regione e un livello; il non noto. Lo stato «libero» è scritto solo per il suolo osservato nel fotogramma corrente; l'avviso di possibile presenza termina solo per osservazione diretta. Nelle ipotesi dichiarate sono dimostrati la sicurezza del «libero», il contenimento del soggetto nascosto nell'insieme di possibile presenza e la condizione sulla fascia di guardia. In simulazione, con modello ideale del sensore, un ostacolo fermo largo $2\unita{m}$, il bordo dell'ombra noto e un soggetto alla volta sul proprio percorso, la regione di previsione è emessa a ogni ciclo dell'occlusione (0 interruzioni su $23\,402$ passaggi) e contiene il soggetto in almeno 9 casi su 10 fino a $2\unita{s}$ e in almeno 85 su 100 fino a $3\unita{s}$. Il classificatore degli oggetti è un componente richiesto, qui simulato e non valutato. Nessuna prova è stata eseguita su un sensore fisico: i risultati valgono a camera ferma, al chiuso, fino a $3\unita{m}$.
